% ============================================================================
%  chapter15.tex
%  The Linux Filesystem Hierarchy
% ============================================================================

\chapter{The Linux Filesystem Hierarchy}

\section{The One Tree}

Every file on a Linux system lives in one tree, hanging from the single
root \code{/}. There is no C:\, no drive letters, no per-device roots.
A USB stick, a network share, and the hard disk all appear as
\emph{parts of the same tree}, attached at directories called
\emph{mount points}. This chapter is a map of that tree: what lives in
each directory, why it lives there, and---more interestingly---how the
layout came to be and what operating-system ideas it encodes.

\begin{verbatim}
   The whole machine, as seen from one path:

        /
        |-- bin/      essential user programs
        |-- sbin/     essential system administration programs
        |-- lib/      shared libraries for /bin and /sbin
        |-- etc/      host-specific configuration
        |-- home/     user home directories
        |-- root/     the superuser's home
        |-- boot/     the kernel and boot images
        |-- dev/      the hardware, as files
        |-- tmp/      world-writable scratch
        |-- var/      variable data: logs, spools, caches
        |-- run/      runtime state since boot
        |-- usr/      the bulk of the system
        |-- opt/      third-party software
        |-- srv/      service data
        |-- mnt/      temporary mount points
        |-- media/    removable media
        |-- proc/     kernel and process info, generated on the fly
        `-- sys/      kernel devices and parameters, on the fly
\end{verbatim}

We have met every idea in this tree already. Chapter~8 gave us the
\emph{namespace}: names pointing to inodes pointing to blocks. This
chapter maps that namespace, adds the history that shaped it, and
explains the OS machinery---mount points, virtual filesystems, device
nodes, tmpfs---that makes one tree possible over many kinds of storage.

\begin{osconcept}
A single tree is an operating system decision with enormous payoff. By
folding every storage device into one namespace, the OS lets programs
refer to anything---a disk, a terminal, a kernel statistic, a network
stream---with one syntax and one set of tools. The cost is that the
tree's \emph{meaning} is not self-evident; a standard layout is needed
so every program can find what it needs. That standard, and the OS
mechanisms behind it, are this chapter's subject.
\end{osconcept}

\section{The Story of the Tree}

The layout was not designed on a whiteboard. It \emph{grew}, one
survival decision at a time, and its quirks are fossils of that
history.

\begin{verbatim}
   1971 -- Version 1 Unix, Bell Labs. The entire operating system runs
   in roughly 32 KB of memory on a single disk. The tree is already
   there, and it is tiny:

        /bin     the handful of programs (cat, ls, sh, cc...)
        /etc     "et cetera" -- configuration and random admin files
        /dev     the devices (a magnetic tape, a disk, a console)
        /usr     the user files -- this is where HOME DIRECTORIES were

   That is the whole story's seed: the first Unix had about five
   directories, and /usr really did mean "user."

   The RK05 disk pack.  The story of /usr's split:
   The root filesystem lived on a small fixed disk. The /usr filesystem
   lived on a larger removable RK05 pack. To boot and repair a broken
   system you could only use the small disk, so the binaries needed to
   recover sat in /bin; everything else -- and the users' files -- went
   to /usr. When the repair disk could not mount /usr, you still had a
   working, minimal system. The split is a backup-and-repair strategy
   fossilized into the layout.
\end{verbatim}

From that seed, two evolutionary pressures shaped the rest. First,
\emph{hardware}: as disks grew, the idea of ``the tiny essential set
vs the rest'' persisted as \code{/bin} vs \code{/usr/bin}, and the
early \code{/usr} directory---which had held home directories---was
colonized by system software, while user files moved to
\code{/home}. Second, \emph{diversity}: BSD and System V rearranged
and disagreed (BSD favored \code{/usr/local} for add-ons, System V
introduced \code{/opt}), and by the early 1990s every Linux
distribution had its own slightly different tree.

\begin{verbatim}
   The tree diverged, and Linux inherited the mess:
     1980s: BSD and System V each reorganize; /home becomes the user
            area; /usr/local, /opt, /var appear in different forms.
     1993:  the Filesystem Hierarchy Standard (FHS) is started to
            make the layout CONSISTENT across Linux distributions --
            so that /etc/passwd means the same thing everywhere, and
            a package built for one distribution finds its libraries
            on another. The FHS (now version 3.0, 2015) is the tree's
            constitution.

   Two modern chapters in the story:
     The "usr merge" (2010s).  Distributions began making /bin,
       /sbin, /lib symlinks to /usr/bin, /usr/sbin, /usr/lib -- one
       copy of the system software, not two. The old split existed for
       repair-on-a-tiny-disk reasons that a 2015 disk no longer has.
     The arrival of /run (2011).  Runtime state (pids, sockets) moved
       out of /var/run into /run, a tmpfs mounted before /var is --
       because early boot needs the state before the big disk is up.
\end{verbatim}

\begin{osconcept}
The tree's history is a history of operating-system constraints. The
root-vs-usr split encodes repair strategy; the FHS encodes
\emph{portability and coordination} (a standard layout is an agreement
that package managers, bootloaders, and administrators all rely on);
the usr-merge and \code{/run} encode the modern boot sequence. Reading
the tree is reading forty years of OS decisions, each one a solution
to a problem of its era.
\end{osconcept}

\section{The OS Concept: Mount Points and the Root Filesystem}

How does one tree hold many disks? The answer is \emph{mounting}. A
\emph{mount point} is a directory in the tree where a filesystem is
\emph{attached}. The kernel hangs the root of one filesystem onto a
directory of another, and from then on the two are one tree. The
\emph{root filesystem} is the one mounted at \code{/}; everything else
is grafted onto it.

\begin{verbatim}
   Two filesystems, one tree:

     before mounting:                      after mounting:
     (root fs on disk A)                   (disk B visible at /home)
       /                                    /
       |-- bin/                             |-- bin/
       |-- etc/                             |-- etc/
       `-- home/  (empty dir,             |-- home/
                 a mount point)                 |-- alice/
             (disk B: /home)                    `-- bob/
       /home                                  (the files of disk B)
       |-- alice/
       `-- bob/

   Mounting is an OS operation (`mount`):
     1. the kernel reads disk B's superblock,
     2. it attaches disk B's root to the directory /home,
     3. from then on, "alice" is found by walking INTO /home.

   Key facts:
     - A mount point is just a directory. Its own contents are
       HIDDEN while a filesystem is mounted over it.
     - Different filesystems can sit at different mounts: ext4 at /,
       tmpfs at /tmp, proc at /proc, a USB drive at /media/usb.
     - Only the kernel can perform the graft, which is why mounting
       is a privileged operation and why /etc/fstab (the table of
       "what mounts where at boot") is a system file.
\end{verbatim}

\begin{lstlisting}[caption={Reading the mount table from inside a program.}]
#include <stdio.h>

int main(void) {
    FILE *f;
    char line[256];

    /* /proc/mounts is the kernel's own list of what is mounted where. */
    f = fopen("/proc/mounts", "r");
    if (f == NULL) {
        perror("fopen");
        return 1;
    }

    while (fgets(line, sizeof(line), f) != NULL) {
        printf("%s", line);     /* device  mount_point  type  options */
    }

    fclose(f);
    return 0;
}
\end{lstlisting}

The output (a few lines of it) looks like:

\begin{verbatim}
   /dev/sda1 / ext4 rw,relatime ...
   proc /proc proc rw,relatime ...
   sysfs /sys sysfs rw,relatime ...
   devtmpfs /dev devtmpfs rw ...
   tmpfs /tmp tmpfs rw ...
   tmpfs /run tmpfs rw ...

   Read it as: "the device (or pseudo-device) on the left is attached
   at the directory in the middle, with the filesystem type on the
   right." / is ext4 on the real disk; /proc, /sys, /dev, /tmp, /run
   are their own (often RAM-backed or virtual) filesystems.
\end{verbatim}

\begin{osconcept}
Mounting is the kernel's namespace surgery: it rewrites the tree's
resolution so that a directory names a different filesystem's root.
Only the kernel can do this, because only the kernel owns the mapping
from paths to inodes and the translation from one filesystem's
structures to another's. The mount table \emph{is} the tree's
scaffolding; the directories you see are the leaves it supports.
\end{osconcept}

\section{The Directories, One by One}

With the machinery in place, here is the tree as a reference. For each
directory: what it holds, what it does, and the problem it solves.

\subsection{The Core}

\begin{verbatim}
   /bin   essential user programs (ls, cp, mv, cat, sh...)
          -- what you need to boot, repair, and use a bare system.
          On modern distributions /bin is a symlink to /usr/bin
          (the "usr merge"): one copy, two names.
          PROBLEM SOLVED: "essential" vs "everything else" was the
          1971 repair strategy; the merge removes the duplicate.

   /sbin  essential ADMINISTRATION programs (mount, fdisk, shutdown...)
          -- the tools that repair and maintain the system itself.
          Also merged into /usr/sbin on modern systems.
          PROBLEM SOLVED: keeping the tools that can break the
          system out of the hands of ordinary PATH lookups.

   /lib   shared libraries for /bin and /sbin (libc, the dynamic
          loader) -- without these, the essential binaries cannot run.
          Merged into /usr/lib.
          PROBLEM SOLVED: one copy of each library, shared by every
          program that uses it (see Chapter 13's memory chapter).

   /etc   host-specific CONFIGURATION. /etc/passwd (users),
          /etc/fstab (what mounts where), /etc/hosts (names),
          /etc/ssh (server keys)... The name is "et cetera": the
          directory where the files that fit nowhere else went.
          PROBLEM SOLVED: separating the machine's own settings from
          the software (which lives in /usr, shared or reinstallable).
          /etc is what an administrator backs up to carry a machine's
          identity to a new one.
\end{verbatim}

\subsection{Users and Data}

\begin{verbatim}
   /home  the user home directories. /home/alice, /home/bob.
          Each user's files, dotfiles, and configuration live here.
          PROBLEM SOLVED: one place for personal data, protected by
          per-user permissions (Chapter 8).

   /root  the superuser's home. (It is not /home/root, so that root's
          home is never hidden by a /home mount or an unmount.)

   /var   VARIABLE data -- content that changes while the machine runs:
          /var/log  (logs, from Chapter 9's daemons),
          /var/spool (queues: mail, print jobs),
          /var/cache (caches that can be rebuilt),
          /var/lib   (stateful data: databases, package state).
          PROBLEM SOLVED: keeping writable, growing data off the
          root filesystem so / (or /usr) can stay read-only or be
          on a smaller, more protected device.

   /tmp   world-writable scratch space. Any process may create files
          here; the sticky bit (Chapter 8) lets the owner delete them.
          Cleared on reboot; often its own tmpfs (RAM-backed).
          PROBLEM SOLVED: a place for throwaway files without a
          per-user allocation.

   /run   runtime state since the last boot: pids of daemons, sockets,
          lock files. A tmpfs -- RAM-backed, cleared at reboot.
          PROBLEM SOLVED: early boot needs state before /var is
          mounted; /run exists from the first moment.

   /srv   service data: the files web and FTP servers actually serve
          (/srv/www). Optional, but when used it is the answer to
          "where does the website live?"
\end{verbatim}

\subsection{The System Proper}

\begin{verbatim}
   /usr   "user system resources" -- the BULK of the operating
          system, and the biggest directory on the disk:
            /usr/bin       most programs
            /usr/sbin      most admin tools
            /usr/lib       most libraries
            /usr/include   C headers
            /usr/share     architecture-independent data (docs, icons,
                           locales, man pages)
            /usr/local     software installed by the ADMIN, by hand
                           (a BSD tradition: the distribution owns
                           /usr; the admin owns /usr/local)
          PROBLEM SOLVED: /usr used to mean "user files"; when home
          directories moved to /home, the name stayed and the system
          software moved in. The merge made /bin a doorway to it.

   /opt   OPTIONAL, third-party software: vendor packages installed as
          their own trees (/opt/google/, /opt/lampp/).
          System V tradition, codified by the FHS so commercial
          software has a place that never collides with the distro's
          files.
          PROBLEM SOLVED: "my package must not fight the distribution."

   /boot  what the machine needs BEFORE the OS runs: the kernel
          (vmlinuz), the initramfs, and the bootloader's files.
          Often a separate, small, early-read partition.
          PROBLEM SOLVED: the bootloader and BIOS read this before
          any filesystem driver exists for the main disk.

   /dev   the DEVICES, as files. /dev/null, /dev/zero, /dev/tty,
          /dev/sda (your disk), /dev/nvme0n1... Writing to a device
          file talks to hardware (Chapter 5's "everything is a file").
          PROBLEM SOLVED: hardware access through the ordinary file
          interface. udev and devtmpfs create the entries at boot as
          devices are discovered.

   /mnt   a convention: temporary mount points (a filesystem you are
          about to inspect, a repair disk). Admins mount things here
          BY HAND.
   /media removable media the desktop mounts automatically (USB
          sticks, CD-ROMs) -- the automatic counterpart of /mnt.
\end{verbatim}

\subsection{Virtual Filesystems: Two Meanings of ``Virtual''}

The word \emph{virtual filesystem} is used for two related but distinct
things in Linux, and keeping them apart prevents a lot of confusion.

\begin{verbatim}
   MEANING 1 -- the VFS layer (the "Virtual File System" switch).
   This is not a filesystem at all; it is an INTERFACE that every
   filesystem must implement. The kernel presents one uniform API --
   open(), read(), write(), stat(), readdir() -- and the VFS hands
   each call to whichever filesystem the path lives on.

        open()/read()/write()/stat()   (the syscalls of ch.5, ch.8)
                        |
                  [ VFS LAYER ]
                  the one interface
             /        |       |       \        \
         ext4      tmpfs    procfs    sysfs    devtmpfs
         (disk)    (RAM)    (kernel)  (kernel)  (kernel)
            each implements the SAME operations, differently

   THIS is the mechanism behind "everything is a file": /dev/null,
   a pipe, an ext4 file, and /proc/meminfo all answer the same
   syscalls because the VFS defines one contract and every
   filesystem fulfills it.

   MEANING 2 -- filesystems with no disk behind them
   (also called PSEUDO-filesystems). Here "virtual" means: the files
   do not exist on any storage device. Their contents are either
   generated on the fly by the kernel (proc, sysfs) or are just RAM
   (tmpfs, devtmpfs). Read /proc/meminfo and the kernel formats its
   memory tables into text at that moment; there is no file to back
   it up. They live only in RAM and vanish at reboot.

   The two meanings meet: pseudo-filesystems (meaning 2) are simply
   filesystems that implement the VFS interface (meaning 1) without
   a storage device.
\end{verbatim}

When Chapter~8 called \code{/proc} a \emph{virtual filesystem}, it was
using meaning 2: files whose bytes are produced on demand. When kernel
documentation talks about \emph{the VFS}, it means meaning 1: the
abstraction layer that makes every kind of filesystem interchangeable.
Both are operating system inventions, and both are why your C program
can treat a kernel statistic and a disk file with the same
\code{open}/\code{read}/\code{close} code.

\begin{osconcept}
The VFS is the OS's plug-in architecture for storage. By defining one
interface and making every filesystem implement it, the kernel lets a
single set of system calls serve disks, RAM, and the kernel's own
internals. Pseudo-filesystems push the idea to its limit: there is no
``storage'' at all, just kernel state rendered as paths. Both exist
because only the kernel can route a syscall to the right implementation
and manufacture ``files'' from its own data structures.
\end{osconcept}

\begin{verbatim}
   /proc  the PROCESS and kernel view, as files -- generated on the
          fly when you read them (Chapter 8). /proc/self/status,
          /proc/cpuinfo, /proc/mounts, /proc/<pid>/... Nothing here
          is stored; it is the kernel's data structures rendered as
          a filesystem.

   /sys   the DEVICE MODEL and kernel parameters, as files. /sys/class,
          /sys/devices, /sys/module... the modern, structured sibling
          of /proc. Writing certain files changes kernel behavior
          ("echo 1 > /sys/.../unload").

   /dev   (listed above) is the third kernel interface: device nodes,
          maintained by udev. All three are VIRTUAL: no disk, no
          persistence, just the file abstraction over kernel state.
\end{verbatim}

\begin{verbatim}
   The five virtual / kernel filesystems, in one table:

     mount     type       backed by        holds
     ------    ---------  ---------------  --------------------------
     /         ext4       the real disk    the tree's home
     /proc     proc       kernel RAM       processes + kernel info
     /sys      sysfs      kernel RAM       the device model
     /dev      devtmpfs   kernel RAM       device nodes
     /tmp, /run, /dev/shm tmpfs           RAM (page-backed)

   "tmpfs" = a filesystem whose storage is RAM pages; it exists while
   the machine runs and vanishes at reboot. That is why /run and /tmp
   are always empty after a clean boot.
\end{verbatim}

\begin{osconcept}
These directories are not decoration; they are the OS's own
bookkeeping made visible. \code{/proc}, \code{/sys}, and \code{/dev}
are the kernel exposing its internals through the file abstraction---
the same design choice as Chapter~8's virtual filesystem, taken to the
whole system. \code{tmpfs} shows the kernel serving RAM as a
filesystem. Reading the tree is reading the kernel.
\end{osconcept}

\section{The Problem the Tree Solves}

With every directory explained, the question ``why does this layout
exist at all?'' has a four-part answer:

\begin{verbatim}
   1. COORDINATION.  A standard layout is an agreement. A program can
      assume /etc/passwd, /usr/include/stdio.h, /var/log exist and
      mean the same thing on every distribution. Without the FHS,
      every package would need a per-distro install script.

   2. SEPARATION OF CONCERNS.  Read-only vs writable (/usr vs /var),
      essential vs optional (/bin vs /opt), static vs stateful
      (/usr/share vs /var/lib), software vs configuration (/usr vs
      /etc). Each split is a different backup, mount, and permission
      policy.

   3. REPAIRABILITY.  The root-vs-usr split's fossil: a minimal,
      bootable, recoverable core exists even when the big filesystem
      is broken. /boot preserves the same idea for the bootloader.

   4. SHARED MACHINERY.  One namespace means one set of tools: find,
      ls, tar, rsync -- and every program in this book -- work on
      disks, devices, and kernel interfaces alike, because they are
      all files.
\end{verbatim}

\begin{ruleofthumb}[Let the layout answer the question for you]
When you wonder where a file belongs, ask \emph{what kind of thing it
is}: configuration goes in \code{/etc}, variable data in
\code{/var}, software in \code{/usr} or \code{/opt}, scratch in
\code{/tmp}. The tree is a decision procedure, and following it keeps
your programs portable and your system repairable.
\end{ruleofthumb}

\section{Practical: A Filesystem Tourist}

Let us build a small tool that walks the root with the directory API
of Chapter~8 and labels each top-level entry with its type---a
machine-readable version of the tree at the top of this chapter.
Everything here is code you already know: \code{opendir},
\code{readdir}, \code{lstat}, and the \code{S\_IS*} macros.

\begin{lstlisting}[caption={fstree.c: walk / and label each top-level entry.}]
#include <stdio.h>
#include <stdlib.h>
#include <string.h>
#include <dirent.h>
#include <sys/stat.h>

/* One-letter label for the file type, as in ls -l (Chapter 8). */
static char type_char(mode_t mode) {
    if (S_ISDIR(mode))  return 'd';
    if (S_ISLNK(mode))  return 'l';
    if (S_ISREG(mode))  return '-';
    if (S_ISCHR(mode))  return 'c';
    if (S_ISBLK(mode))  return 'b';
    if (S_ISFIFO(mode)) return 'p';
    if (S_ISSOCK(mode)) return 's';
    return '?';
}

/* Print the type and name of every entry in one directory. */
static void tour(const char *path) {
    DIR *d;
    struct dirent *e;

    d = opendir(path);
    if (d == NULL) {
        perror("opendir");
        return;
    }

    while (1) {
        struct stat st;
        char full[1024];
        int plen;

        e = readdir(d);
        if (e == NULL) {
            break;
        }
        if (e->d_name[0] == '.') {
            continue;          /* skip . and .. */
        }

        plen = snprintf(full, sizeof(full), "%s/%s", path, e->d_name);
        if (plen < 0 || plen >= (int) sizeof(full)) {
            continue;
        }

        if (lstat(full, &st) < 0) {
            continue;
        }

        printf("%c %s\n", type_char(st.st_mode), e->d_name);
    }

    closedir(d);
}

int main(int argc, char *argv[]) {
    const char *root;

    root = (argc >= 2) ? argv[1] : "/";
    printf("type  entry  (under %s)\n", root);
    tour(root);
    return 0;
}
\end{lstlisting}

\begin{verbatim}
   $ ./fstree
   type  entry  (under /)
   d bin
   d boot
   d dev
   d etc
   d home
   d lib
   d media
   d mnt
   d opt
   d proc
   d root
   d run
   d sbin
   d srv
   d sys
   d tmp
   d usr
   d var

   On a system with the usr-merge, ls -l /bin reveals the history at a
   glance: l bin -> usr/bin. The 1971 split has become a symlink.
\end{verbatim}

\begin{researchnote}[The tree's constitution]
The Filesystem Hierarchy Standard is the document this chapter has
been quoting, and it is worth reading from the source. The current
edition, FHS 3.0, published by the Linux Foundation in 2015, specifies
the layout precisely enough that a distribution can be checked against
it mechanically---which is exactly what makes cross-distribution
software possible.

\begin{verbatim}
   The FHS's stated purpose, condensed from its own preface:

       "The filesystem hierarchy standard describes a conventional
        layout that allows software to predict where its files and
        data reside, so that software can be written to work on any
        conforming system."

   Every table in this chapter--/bin vs /usr/bin, /etc for
   configuration, /var for variable data--is the FHS made flesh. It is
   a rare thing: a standard that is simultaneously bureaucratic and
   load-bearing.
\end{verbatim}

Reading the FHS straight through is the fastest way to cement this
chapter: it is short, unambiguous, and each entry is the "problem
solved" text of Section 15.4 in its most careful form.

\paperref{Filesystem Hierarchy Standard, Version 3.0, Linux Foundation,
2015}{https://refspecs.linuxfoundation.org/FHS_3.0/fhs-3.0.pdf}
\end{researchnote}

\section{Structures in Bytes}

The hierarchy's byte structures are mostly on disk. A symbolic link is
the simplest: its data bytes are the target path.

\begin{verbatim}
   symlink "data_soft.txt" -> "data.txt"

   inode (mode = symlink)      data
   +-----------------+         +-------------+
   | size = 8        |-------->| "data.txt\0" |   8 bytes + NUL
   +-----------------+         +-------------+
\end{verbatim}

The ext4 superblock (1024 bytes, at offset 1024 of the partition)
holds the fields the kernel reads to mount a filesystem:

\begin{verbatim}
   ext4 superblock, key offsets:

   offset  size  field
   0       4     s_inodes_count
   4       4     s_blocks_count
   20      4     s_log_block_size    (block size = 1024 << this)
   56      2     s_magic             (0xEF53 marks a real ext4)
   88      2     s_inode_size        (256 on ext4)
\end{verbatim}

And \code{/proc/mounts} renders the kernel's mount table as text---one
line per mount, each field separated by spaces:

\begin{verbatim}
   /dev/sda1 / ext4 rw,relatime 0 0
   ^device   ^mount point ^fstype ^options ^dump ^pass
\end{verbatim}

Everything the tree is built from---inodes, superblocks, mount
records---is bytes; the hierarchy is the names that point at them.

\section{Summary}

\begin{verbatim}
   The filesystem hierarchy, in one paragraph:

   one tree from / ; every device, device-node, and kernel interface
   hangs on it at a MOUNT POINT; the layout is a 1971 fossil (bin vs
   usr, /etc) standardized by the FHS (coordination, separation,
   repairability) and modernized by the usr-merge and /run.

   The directories as a decision procedure:
     configuration  /etc      variable data  /var
     software       /usr,/opt scratch        /tmp
     users          /home     admin          /sbin,/root
     boot           /boot     services       /srv
     hardware       /dev      kernel         /proc,/sys
     mounting       /mnt,/media

   Rules that keep you oriented:
     1. ls -l /bin is worth doing once on your own machine: the
        history is visible in the symlinks;
     2. /proc/mounts tells you what the tree is built from;
     3. when in doubt about where a file belongs, follow the FHS's
        "what kind of thing is it?" question.
\end{verbatim}

\begin{osconcept}
The hierarchy is the OS's address book made public. It works only
because the kernel owns mounting, only because virtual filesystems let
the kernel render itself as paths, and only because a standard keeps
the agreement stable across distributions. This chapter sat between
the code you have written and the machine it runs on: the tree is
where every path in every earlier chapter actually lived.
\end{osconcept}

The tour is complete. The next chapter returns to writing code and
builds the machine's other great promise: durable storage. We design a
wire format, append and replay a log, and make data survive a
crash---before the final chapter carries it across the network.

